\subsection{Symmetric and Exterior Powers}

% This subsection develops symmetric and exterior powers from multilinear maps
% and tensor powers, then assembles the associated graded algebras.

\begin{definition} \label{def:special-multilinear-maps} \hlt{Symmetric, Skew-Symmetric, and Alternating Maps}\\
Let $n\in\mathbb{Z}_{>0}$, let $\mathcal{V},\mathcal{W}$ be $\mathbb{F}$-vector spaces, and set $P_n\defeq\{(i,j)\in\{1,\ldots,n\}^2:i<j\}$. Let $f\in\operatorname{Mult}_{\mathbb{F}}(\mathcal{V},\ldots,\mathcal{V};\mathcal{W})$, with $n$ copies of $\mathcal{V}$.\\
For $\sigma\in\mathfrak{S}_n$, set $f^\sigma(\mathbf{v}_1,\ldots,\mathbf{v}_n)\defeq f(\mathbf{v}_{\sigma(1)},\ldots,\mathbf{v}_{\sigma(n)})$. For $(i,j)\in P_n$, let $s_{ij}$ transpose $i$ and $j$.
\begin{enumerate}[label=(\roman*)]
\setlength{\itemsep}{0pt}
\item The map $f$ is \hlt{symmetric} if $(\forall\,(i,j)\in P_n)\,f^{s_{ij}}=f$.
\item The map $f$ is \hlt{skew-symmetric} if $(\forall\,(i,j)\in P_n)\,f^{s_{ij}}=-f$.
\item The map $f$ is \hlt{alternating} if $(\forall\,\mathbf{v}\in\mathcal{V}^{\times n})[(\exists\,(i,j)\in P_n)\,\mathbf{v}_i=\mathbf{v}_j\Rightarrow f(\mathbf{v})=0]$.
\end{enumerate}
\end{definition}

\begin{proposition} \label{prop:alternating-skew-characteristic} \hlt{Alternation, Skew-Symmetry, and Characteristic}\\
Let $n\in\mathbb{Z}_{>0}$, let $\mathcal{V},\mathcal{W}$ be $\mathbb{F}$-vector spaces, and set $P_n\defeq\{(i,j)\in\{1,\ldots,n\}^2:i<j\}$. Let $f\in\operatorname{Mult}_{\mathbb{F}}(\mathcal{V},\ldots,\mathcal{V};\mathcal{W})$, with $n$ copies of $\mathcal{V}$.
\begin{enumerate}[label=(\roman*)]
\setlength{\itemsep}{0pt}
\item If $f$ is alternating, then $f$ is skew-symmetric.
\item If $\operatorname{char}(\mathbb{F})\neq2$ and $f$ is skew-symmetric, then $f$ is alternating.
\item If $\operatorname{char}(\mathbb{F})=2$, then the symmetric and skew-symmetric conditions coincide.
\end{enumerate}
\begin{proof}
\begin{enumerate}[label=(\roman*)]
\setlength{\itemsep}{0pt}
\item Fix $(i,j)\in P_n$ and the entries outside slots $i,j$, and let $g:\mathcal{V}^2\to\mathcal{W}$ be the resulting bilinear map.\\
Then $(\forall\,\mathbf{u}\in\mathcal{V})\,g(\mathbf{u},\mathbf{u})=0$ by alternation. For $\mathbf{u},\mathbf{v}\in\mathcal{V}$, $0=g(\mathbf{u}+\mathbf{v},\mathbf{u}+\mathbf{v})=g(\mathbf{u},\mathbf{v})+g(\mathbf{v},\mathbf{u})$.\\
Hence $g(\mathbf{v},\mathbf{u})=-g(\mathbf{u},\mathbf{v})$. The fixed data were arbitrary, so $(\forall\,(i,j)\in P_n)\,f^{s_{ij}}=-f$.
\item Fix $\mathbf{v}\in\mathcal{V}^{\times n}$ and $(i,j)\in P_n$ with $\mathbf{v}_i=\mathbf{v}_j$. Then $f(\mathbf{v})=f^{s_{ij}}(\mathbf{v})=-f(\mathbf{v})$ by skew-symmetry, whence $2f(\mathbf{v})=0$.\\
The scalar $2$ is invertible in $\mathbb{F}$ by $\operatorname{char}(\mathbb{F})\neq2$, so $f(\mathbf{v})=0$.
\item The equality $-1=1$ holds in $\mathbb{F}$ by $\operatorname{char}(\mathbb{F})=2$, so the two conditions coincide.
\end{enumerate}
\hlt{Structural Commentary}: Diagonal polarization establishes skew-symmetry; invertibility of $2$ recovers alternation.\\
\hlt{Key Insight}: Characteristic $2$ removes the sign distinction but not diagonal vanishing.
\end{proof}
\end{proposition}

\begin{definition} \label{def:permutation-action-on-tensor-powers} \hlt{Permutation Action on Tensor Powers}\\
Let $\mathcal{V}$ be an $\mathbb{F}$-vector space and $p\in\mathbb{Z}_{\geq0}$. Set $T_0\defeq\mathbb{F}$, $T_p\defeq\mathcal{V}^{\otimes p}$ for $p>0$, and $G_p\defeq\mathfrak{S}_p$.\\
For $p=0$, let $t_0:\{()\}\to T_0$ satisfy $t_0(())=1$ and set $\lambda_{\operatorname{id}}\defeq I_{T_0}$.\\
For $p>0$, let $t_p:\mathcal{V}^{\times p}\to T_p$ be canonical and set $\lambda_{\sigma}\defeq\tau_{\sigma^{-1}}\in\LL(T_p)$ for $\sigma\in G_p$ by Theorem \ref{thm:iterated-tensor-products-basic-laws}(iv).\\
Thus $\lambda_{\sigma}t_p(\mathbf{v}_1,\ldots,\mathbf{v}_p)=t_p(\mathbf{v}_{\sigma^{-1}(1)},\ldots,\mathbf{v}_{\sigma^{-1}(p)})$.
\end{definition}

\begin{proposition} \label{prop:permutation-representation-on-tensor-powers} \hlt{Permutation Representation on Tensor Powers}\\
Let $\mathcal{V}$ be an $\mathbb{F}$-vector space and $p\in\mathbb{Z}_{\geq0}$. Let $T_p$, $G_p$, and $(\lambda_{\sigma})_{\sigma\in G_p}$ be as in Definition \ref{def:permutation-action-on-tensor-powers}.
\begin{enumerate}[label=(\roman*)]
\setlength{\itemsep}{0pt}
\item $\lambda_{\operatorname{id}}=I_{T_p}$.
\item $(\forall\,\sigma,\tau\in G_p)\,\lambda_{\sigma}\lambda_{\tau}=\lambda_{\sigma\tau}$.
\end{enumerate}
\begin{proof}
Both claims hold by Definition \ref{def:permutation-action-on-tensor-powers} when $p=0$. Assume $p>0$.
\begin{enumerate}[label=(\roman*)]
\setlength{\itemsep}{0pt}
\item $(\forall\,\mathbf{v}\in\mathcal{V}^{\times p})\,\lambda_{\operatorname{id}}t_p(\mathbf{v})=t_p(\mathbf{v})$. Hence $\lambda_{\operatorname{id}}=I_{T_p}$ by Theorem \ref{thm:iterated-tensor-products-basic-laws}(i)(b).
\item For $(\sigma,\tau,\mathbf{v})\in G_p^2\times\mathcal{V}^{\times p}$, $\lambda_{\sigma}\lambda_{\tau}t_p(\mathbf{v})=t_p(\mathbf{v}_{\tau^{-1}\sigma^{-1}(1)},\ldots,\mathbf{v}_{\tau^{-1}\sigma^{-1}(p)})=\lambda_{\sigma\tau}t_p(\mathbf{v})$.\\
Thus $(\forall\,\sigma,\tau\in G_p)\,\lambda_{\sigma}\lambda_{\tau}=\lambda_{\sigma\tau}$ by Theorem \ref{thm:iterated-tensor-products-basic-laws}(i)(b).
\end{enumerate}
\hlt{Structural Commentary}: Decomposable evaluation verifies each law; spanning extends it to $T_p$.\\
\hlt{Key Insight}: Inverse indexing converts factor permutations into a left action.
\end{proof}
\end{proposition}

\begin{definition} \label{def:symmetric-and-exterior-tensor-spaces} \hlt{Symmetric and Exterior Powers}\\
Let $\mathcal{V}$ be an $\mathbb{F}$-vector space, let $p\in\mathbb{Z}_{\geq0}$, and use the notation of Definition \ref{def:permutation-action-on-tensor-powers}.\\
Set $P_p\defeq\{(i,j)\in\{1,\ldots,p\}^2:i<j\}$ and $\Delta_p(\mathcal{V})\defeq\{\mathbf{v}\in\mathcal{V}^{\times p}:(\exists\,(i,j)\in P_p)\,\mathbf{v}_i=\mathbf{v}_j\}$.
\begin{enumerate}[label=(\roman*)]
\setlength{\itemsep}{0pt}
\item The \hlt{symmetric tensors} are $\operatorname{ST}^{p}(\mathcal{V})\defeq\{\mathbf{z}\in T_p:(\forall\,\sigma\in G_p)\,\lambda_{\sigma}\mathbf{z}=\mathbf{z}\}$.
\item The \hlt{antisymmetric tensors} are $\operatorname{AT}^{p}(\mathcal{V})\defeq\{\mathbf{z}\in T_p:(\forall\,\sigma\in G_p)\,\lambda_{\sigma}\mathbf{z}=\sign(\sigma)\mathbf{z}\}$.
\item Set $R_{\vee,p}\defeq\operatorname{span}\{\lambda_{\sigma}\mathbf{z}-\mathbf{z}:\sigma\in G_p,\ \mathbf{z}\in T_p\}$.
\item Set $R_{\wedge,p}\defeq\operatorname{span}\{t_p(\mathbf{v}):\mathbf{v}\in\Delta_p(\mathcal{V})\}$.
\item The \hlt{symmetric power} is $\operatorname{Sym}^{p}(\mathcal{V})\defeq T_p/R_{\vee,p}$. For the quotient map $q_{\vee,p}$, set $q_{\vee,p}t_p(\mathbf{v})\defeq\mathbf{v}_1\vee\cdots\vee\mathbf{v}_p$.
\item The \hlt{exterior power} is $\bigwedge^{p}\mathcal{V}\defeq T_p/R_{\wedge,p}$. For the quotient map $q_{\wedge,p}$, set $q_{\wedge,p}t_p(\mathbf{v})\defeq\mathbf{v}_1\wedge\cdots\wedge\mathbf{v}_p$.
\end{enumerate}
For $p=0$, each empty product denotes $1\in\mathbb{F}$.
\end{definition}

\begin{theorem} \label{thm:symmetric-exterior-universal-properties} \hlt{Universal Properties of Symmetric and Exterior Powers}\\
Let $\mathcal{V},\mathcal{U}$ be $\mathbb{F}$-vector spaces and $p\in\mathbb{Z}_{>0}$. Set $M\defeq\operatorname{Mult}_{\mathbb{F}}(\mathcal{V},\ldots,\mathcal{V};\mathcal{U})$, with $p$ copies of $\mathcal{V}$.\\
Let $t_p:\mathcal{V}^{\times p}\to T_p$ be canonical. For $f\in M$ and $\sigma\in G_p$, set $f^\sigma(\mathbf{v}_1,\ldots,\mathbf{v}_p)\defeq f(\mathbf{v}_{\sigma(1)},\ldots,\mathbf{v}_{\sigma(p)})$.
\begin{enumerate}[label=(\roman*)]
\setlength{\itemsep}{0pt}
\item $M_{\vee}\defeq\{f\in M:(\forall\,\sigma\in G_p)\,f^\sigma=f\}$. Map $\Phi_{\vee}:\LL(\operatorname{Sym}^{p}(\mathcal{V}),\mathcal{U})\to M_{\vee}$ defined by $\Phi_{\vee}(T)\defeq T\circ q_{\vee,p}\circ t_p$ is linear isomorphism.
\item $M_{\wedge}\defeq\{f\in M:(\forall\,\mathbf{v}\in\Delta_p(\mathcal{V}))\,f(\mathbf{v})=0\}$. Map $\Phi_{\wedge}:\LL(\bigwedge^{p}\mathcal{V},\mathcal{U})\to M_{\wedge}$ defined by $\Phi_{\wedge}(T)\defeq T\circ q_{\wedge,p}\circ t_p$ is linear isomorphism.
\end{enumerate}
\begin{proof}
\begin{enumerate}[label=(\roman*)]
\setlength{\itemsep}{0pt}
\item For $T\in\LL(\operatorname{Sym}^{p}(\mathcal{V}),\mathcal{U})$, $(\forall\,\sigma\in G_p)\,q_{\vee,p}\lambda_{\sigma}=q_{\vee,p}$ by the definition of $R_{\vee,p}$. Hence $\Phi_{\vee}(T)\in M_{\vee}$.\\
Map $\Phi_{\vee}$ is linear by pointwise operations, linear composition. For $f\in M_{\vee}$, $(\exists!\,\widetilde{f}\in\LL(T_p,\mathcal{U}))\,f=\widetilde{f}\circ t_p$ by Theorem \ref{thm:iterated-tensor-products-basic-laws}(ii).\\
The identity $(\forall\,\sigma\in G_p)\,\widetilde{f}\lambda_{\sigma}=\widetilde{f}$ holds on decomposable tensors by symmetry and on $T_p$ by Theorem \ref{thm:iterated-tensor-products-basic-laws}(i)(b).\\
Hence $R_{\vee,p}\leq\ker(\widetilde{f})$. Thus $(\exists!\,\overline{f}\in\LL(\operatorname{Sym}^{p}(\mathcal{V}),\mathcal{U}))\,\widetilde{f}=\overline{f}\circ q_{\vee,p}$ by Theorem \ref{thm:quotient-universal-property}(i).\\
For $T\in\LL(\operatorname{Sym}^{p}(\mathcal{V}),\mathcal{U})$, the construction applied to $\Phi_{\vee}(T)$ has $\widetilde{\Phi_{\vee}(T)}=T\circ q_{\vee,p}$. Thus $\overline{\Phi_{\vee}(T)}=T$ by uniqueness.\\
For $f\in M_{\vee}$, $\Phi_{\vee}(\overline{f})=\overline{f}\circ q_{\vee,p}\circ t_p=\widetilde{f}\circ t_p=f$. Thus $\Phi_{\vee}$ is an isomorphism.
\item For $T\in\LL(\bigwedge^{p}\mathcal{V},\mathcal{U})$, $\Phi_{\wedge}(T)\in M_{\wedge}$ by the definition of $R_{\wedge,p}$. The map $\Phi_{\wedge}$ is linear by pointwise operations and linear composition.\\
For $f\in M_{\wedge}$, $(\exists!\,\widetilde{f}\in\LL(T_p,\mathcal{U}))\,f=\widetilde{f}\circ t_p$ by Theorem \ref{thm:iterated-tensor-products-basic-laws}(ii).\\
The inclusion $R_{\wedge,p}\leq\ker(\widetilde{f})$ holds by alternation. Thus $(\exists!\,\overline{f}\in\LL(\bigwedge^{p}\mathcal{V},\mathcal{U}))\,\widetilde{f}=\overline{f}\circ q_{\wedge,p}$ by Theorem \ref{thm:quotient-universal-property}(i).\\
For $T\in\LL(\bigwedge^{p}\mathcal{V},\mathcal{U})$, the construction applied to $\Phi_{\wedge}(T)$ has $\widetilde{\Phi_{\wedge}(T)}=T\circ q_{\wedge,p}$. Thus $\overline{\Phi_{\wedge}(T)}=T$ by uniqueness.\\
For $f\in M_{\wedge}$, $\Phi_{\wedge}(\overline{f})=\overline{f}\circ q_{\wedge,p}\circ t_p=\widetilde{f}\circ t_p=f$. Thus $\Phi_{\wedge}$ is an isomorphism.
\end{enumerate}
\hlt{Structural Commentary}: Tensor universality linearizes each map; quotient descent removes its relation subspace.\\
\hlt{Key Insight}: Kernel containment is the exact descent criterion for both tensor quotients.
\end{proof}
\end{theorem}

\begin{theorem} \label{thm:symmetric-exterior-dimensions} \hlt{Dimensions of Symmetric and Exterior Powers}\\
Let $\dim\mathcal{V}=n\in\mathbb{Z}_{>0}$, let $p\in\mathbb{Z}_{\geq0}$, and let $E=(\mathbf{e}_1,\ldots,\mathbf{e}_n)$ be an ordered basis. If $p=0$, set $I_{\vee}=I_{\wedge}\defeq\{()\}$.\\
If $p>0$, set $I_{\vee}\defeq\{(i_1,\ldots,i_p)\in\{1,\ldots,n\}^{p}:i_1\leq\cdots\leq i_p\}$ and $I_{\wedge}\defeq\{(i_1,\ldots,i_p)\in\{1,\ldots,n\}^{p}:i_1<\cdots<i_p\}$.\\
For $I=(i_1,\ldots,i_p)$, set $\mathbf{e}_{I}^{\vee}\defeq\mathbf{e}_{i_1}\vee\cdots\vee\mathbf{e}_{i_p}$ and $\mathbf{e}_{I}^{\wedge}\defeq\mathbf{e}_{i_1}\wedge\cdots\wedge\mathbf{e}_{i_p}$, with empty products equal to $1\in\mathbb{F}$.\\
Set $B_{\vee}\defeq\{\mathbf{e}_{I}^{\vee}:I\in I_{\vee}\}$ and $B_{\wedge}\defeq\{\mathbf{e}_{I}^{\wedge}:I\in I_{\wedge}\}$.
\begin{enumerate}[label=(\roman*)]
\setlength{\itemsep}{0pt}
\item The following symmetric-power conclusions hold.
\begin{enumerate}[label=(\alph*)]
\setlength{\itemsep}{0pt}
\item The set $B_{\vee}$ is a basis of $\operatorname{Sym}^{p}(\mathcal{V})$.
\item $\dim\operatorname{Sym}^{p}(\mathcal{V})=\#I_{\vee}=\binom{n+p-1}{p}$.
\end{enumerate}
\item The following exterior-power conclusions hold, with $\binom{n}{p}\defeq0$ for $p>n$.
\begin{enumerate}[label=(\alph*)]
\setlength{\itemsep}{0pt}
\item The set $B_{\wedge}$ is a basis of $\bigwedge^{p}\mathcal{V}$.
\item $\dim\bigwedge^{p}\mathcal{V}=\#I_{\wedge}=\binom{n}{p}$.
\end{enumerate}
\end{enumerate}
\begin{proof}
For $p>0$, let $t_p:\mathcal{V}^{\times p}\to T_p$ be canonical and, for finite $S$, set $\mathcal{P}_p(S)\defeq\{A\subseteq S:\#A=p\}$.\\
For $p>0$, the tensor words $\mathbf{e}_{j_1}\otimes\cdots\otimes\mathbf{e}_{j_p}$ form a basis of $T_p$ by repeated applications of Theorem \ref{thm:tensor-bases-dimension-and-rank}(i).
\begin{enumerate}[label=(\roman*)]
\setlength{\itemsep}{0pt}
\item
\begin{enumerate}[label=(\alph*)]
\setlength{\itemsep}{0pt}
\item If $p=0$, then $R_{\vee,0}=\{0\}$ and $B_{\vee}=\{1\}$, so $B_{\vee}$ is a basis.\\
Assume $p>0$. Each tensor word is congruent modulo $R_{\vee,p}$ to its nondecreasing rearrangement; hence $B_{\vee}$ spans.\\
For $I\in I_{\vee}$, define $\chi_I\in T_p'$ on the tensor-word basis by $\chi_I(\mathbf{e}_{j_1}\otimes\cdots\otimes\mathbf{e}_{j_p})=1$ when $(j_1,\ldots,j_p)$ is a rearrangement of $I$, and $0$ otherwise. The inclusion $R_{\vee,p}\leq\ker(\chi_I)$ holds by $(\forall\,\sigma\in G_p)\,\chi_I\lambda_{\sigma}=\chi_I$.\\
Thus $(\exists!\,\overline{\chi}_I\in(\operatorname{Sym}^{p}(\mathcal{V}))')\,\chi_I=\overline{\chi}_I\circ q_{\vee,p}$ by Theorem \ref{thm:quotient-universal-property}(i).\\
For $I,J\in I_{\vee}$, $\overline{\chi}_I(\mathbf{e}_{J}^{\vee})=\delta_{I,J}$; hence $B_{\vee}$ is linearly independent and therefore a basis.
\item If $p=0$, then $\dim\operatorname{Sym}^{0}(\mathcal{V})=\#I_{\vee}=1=\binom{n-1}{0}$ by Proposition \ref{prop:binomial-coefficient-edge-and-symmetry}(i).\\
For $p>0$, the map $(i_1,\ldots,i_p)\mapsto\{i_1,i_2+1,\ldots,i_p+p-1\}$ is a bijection from $I_{\vee}$ to $\mathcal{P}_p(\{1,\ldots,n+p-1\})$.\\
Inverse sends $\{a_1<\cdots<a_p\}$ to $(a_1,a_2-1,\ldots,a_p-p+1)$. Thus $\dim\operatorname{Sym}^{p}(\mathcal{V})=\#I_{\vee}=\binom{n+p-1}{p}$ by (a), Proposition \ref{prop:unordered-samples-binomial-coefficients}.
\end{enumerate}
\item
\begin{enumerate}[label=(\alph*)]
\setlength{\itemsep}{0pt}
\item If $p=0$, then $R_{\wedge,0}=\{0\}$ and $B_{\wedge}=\{1\}$, so $B_{\wedge}$ is a basis.\\
Assume $p>0$. If $p\geq2$, fix two slots and all other factors, and write $w(\mathbf{u},\mathbf{v})$ for the resulting tensor word.\\
Then $w(\mathbf{u},\mathbf{v})+w(\mathbf{v},\mathbf{u})=w(\mathbf{u}+\mathbf{v},\mathbf{u}+\mathbf{v})-w(\mathbf{u},\mathbf{u})-w(\mathbf{v},\mathbf{v})\in R_{\wedge,p}$. Thus a transposition changes sign modulo $R_{\wedge,p}$.\\
Repeated basis words lie in $R_{\wedge,p}$ by definition, and every other basis word is congruent up to sign to an increasing word; hence $B_{\wedge}$ spans. Let $E'=(\varepsilon_1,\ldots,\varepsilon_n)$ be the dual basis. For $I=(i_1,\ldots,i_p)\in I_{\wedge}$, set $a_I(\mathbf{v}_1,\ldots,\mathbf{v}_p)\defeq\det[\varepsilon_{i_r}(\mathbf{v}_s)]_{r,s=1}^{p}$.\\
The form $a_I$ is alternating by Theorem \ref{thm:determinant-basic-laws}(ii)--(iv). Thus $(\exists!\,\ell_I\in(\bigwedge^{p}\mathcal{V})')\,a_I=\ell_I\circ q_{\wedge,p}\circ t_p$ by Theorem \ref{thm:symmetric-exterior-universal-properties}(ii).\\
For $I,J\in I_{\wedge}$, $\ell_I(\mathbf{e}_{J}^{\wedge})=\delta_{I,J}$; hence $B_{\wedge}$ is linearly independent and therefore a basis.
\item If $p=0$, then $\dim\bigwedge^{0}\mathcal{V}=\#I_{\wedge}=1=\binom{n}{0}$ by Proposition \ref{prop:binomial-coefficient-edge-and-symmetry}(i).\\
For $0<p\leq n$, the map $(i_1,\ldots,i_p)\mapsto\{i_1,\ldots,i_p\}$ is a bijection from $I_{\wedge}$ to $\mathcal{P}_p(\{1,\ldots,n\})$.\\
Thus $\dim\bigwedge^{p}\mathcal{V}=\#I_{\wedge}=\binom{n}{p}$ by (a) and Proposition \ref{prop:unordered-samples-binomial-coefficients}.\\
If $p>n$, then $I_{\wedge}=\emptyset$ and $\dim\bigwedge^{p}\mathcal{V}=0=\binom{n}{p}$ by (a) and the stated convention.
\end{enumerate}
\end{enumerate}
\hlt{Structural Commentary}: Orbit-constant functionals separate symmetric words; determinants separate increasing exterior words.\\
\hlt{Key Insight}: Dual separators turn spanning normal forms into bases.
\end{proof}
\end{theorem}

\begin{proposition} \label{prop:symmetrization-map} \hlt{Symmetrization and Antisymmetrization in Characteristic Zero}\\
Let $\mathcal{V}$ be $\mathbb{F}$-vector space, $p\in\mathbb{Z}_{\geq0}$, $\operatorname{char}(\mathbb{F})=0$. The scalar $p!$ is nonzero in $\mathbb{F}$.\\
Define $S\defeq\frac{1}{p!}\sum\limits_{\sigma\in G_p}\lambda_{\sigma}\in\LL(T_p)$; $A\defeq\frac{1}{p!}\sum\limits_{\sigma\in G_p}\sign(\sigma)\lambda_{\sigma}\in\LL(T_p)$.
\begin{enumerate}[label=(\roman*)]
\setlength{\itemsep}{0pt}
\item $\operatorname{im}S=\operatorname{ST}^{p}(\mathcal{V})$.
\item $S^2=S$.
\item $\operatorname{im}A=\operatorname{AT}^{p}(\mathcal{V})$.
\item $A^2=A$.
\end{enumerate}
\begin{proof}
\begin{enumerate}[label=(\roman*)]
\setlength{\itemsep}{0pt}
\item For $\tau\in G_p$, $\lambda_{\tau}S=\frac{1}{p!}\sum\limits_{\sigma\in G_p}\lambda_{\tau\sigma}=\frac{1}{p!}\sum\limits_{\rho\in G_p}\lambda_{\rho}=S$ by the reindexing $\rho=\tau\sigma$. Hence $\operatorname{im}S\subseteq\operatorname{ST}^{p}(\mathcal{V})$.\\
If $\mathbf{z}\in\operatorname{ST}^{p}(\mathcal{V})$, then $S\mathbf{z}=\frac{1}{p!}\sum\limits_{\sigma\in G_p}\mathbf{z}=\mathbf{z}$ by Proposition \ref{prop:ordered-sampling-with-and-without-replacement}(iii). Thus $\operatorname{ST}^{p}(\mathcal{V})\subseteq\operatorname{im}S$; $\operatorname{im}S=\operatorname{ST}^{p}(\mathcal{V})$.
\item For $\mathbf{x}\in T_p$, $S\mathbf{x}\in\operatorname{ST}^{p}(\mathcal{V})$ by (i), so $S^2\mathbf{x}=S\mathbf{x}$ by the defining average and Proposition \ref{prop:ordered-sampling-with-and-without-replacement}(iii). Hence $S^2=S$.
\item For $p>0$, let $(\mathbf{e}_1,\ldots,\mathbf{e}_p)$ be the standard basis of $\mathbb{F}^{p}$, and let $P_{\sigma}$ have columns $\mathbf{e}_{\sigma(1)},\ldots,\mathbf{e}_{\sigma(p)}$ for $\sigma\in G_p$.\\
Then $\det P_{\sigma}=\sign(\sigma)$ by Definition \ref{def:matrix-trace-and-determinant}(ii), and $P_{\sigma}P_{\tau}=P_{\sigma\tau}$.\\
Hence $\sign(\sigma\tau)=\sign(\sigma)\sign(\tau)$ by Theorem \ref{thm:determinant-basic-laws}(vi). The identity also holds for $p=0$.\\
For $\tau\in G_p$, $\lambda_{\tau}A=\frac{1}{p!}\sum\limits_{\sigma\in G_p}\sign(\sigma)\lambda_{\tau\sigma}=\sign(\tau)A$ by the reindexing $\rho=\tau\sigma$. Hence $\operatorname{im}A\subseteq\operatorname{AT}^{p}(\mathcal{V})$.\\
If $\mathbf{z}\in\operatorname{AT}^{p}(\mathcal{V})$, then $A\mathbf{z}=\frac{1}{p!}\sum\limits_{\sigma\in G_p}\sign(\sigma)^2\mathbf{z}=\mathbf{z}$ by Proposition \ref{prop:ordered-sampling-with-and-without-replacement}(iii). Thus $\operatorname{AT}^{p}(\mathcal{V})\subseteq\operatorname{im}A$; $\operatorname{im}A=\operatorname{AT}^{p}(\mathcal{V})$.
\item For $\mathbf{x}\in T_p$, $A\mathbf{x}\in\operatorname{AT}^{p}(\mathcal{V})$ by (iii), so $A^2\mathbf{x}=A\mathbf{x}$ by the defining average and Proposition \ref{prop:ordered-sampling-with-and-without-replacement}(iii). Hence $A^2=A$.
\end{enumerate}
\hlt{Structural Commentary}: Reindexed character averages identify the images; fixed-point identities establish idempotence.\\
\hlt{Key Insight}: Normalisation by $p!$ turns each character average into a projection.
\end{proof}
\end{proposition}

\begin{definition} \label{def:graded-tensor-symmetric-exterior-algebras} \hlt{Tensor, Symmetric, and Exterior Algebras}\\
Let $\mathcal{V}$ be an $\mathbb{F}$-vector space.
\begin{enumerate}[label=(\roman*)]
\setlength{\itemsep}{0pt}
\item The \hlt{tensor algebra} is $\mathcal{T}(\mathcal{V})\defeq\bigoplus\limits_{p\geq0}T_p$. Unit is $1\in T_0=\mathbb{F}$, homogeneous multiplication is tensor concatenation.
\item The \hlt{symmetric algebra} is $\operatorname{Sym}(\mathcal{V})\defeq\bigoplus\limits_{p\geq0}\operatorname{Sym}^{p}(\mathcal{V})$. For $\mathbf{x}\in T_p$ and $\mathbf{y}\in T_q$, set $q_{\vee,p}(\mathbf{x})\vee q_{\vee,q}(\mathbf{y})\defeq q_{\vee,p+q}(\mathbf{x}\otimes\mathbf{y})$.
\item The \hlt{exterior algebra} is $\bigwedge\mathcal{V}\defeq\bigoplus\limits_{p\geq0}\bigwedge^{p}\mathcal{V}$. For $\mathbf{x}\in T_p$ and $\mathbf{y}\in T_q$, set $q_{\wedge,p}(\mathbf{x})\wedge q_{\wedge,q}(\mathbf{y})\defeq q_{\wedge,p+q}(\mathbf{x}\otimes\mathbf{y})$.
\end{enumerate}
\end{definition}

\begin{proposition} \label{prop:graded-tensor-symmetric-exterior-algebra-laws} \hlt{Graded Algebra Laws}\\
Let $\mathcal{V}$ be an $\mathbb{F}$-vector space, with the three graded constructions and products from Definition \ref{def:graded-tensor-symmetric-exterior-algebras}.
\begin{enumerate}[label=(\roman*)]
\setlength{\itemsep}{0pt}
\item Tensor concatenation makes $\mathcal{T}(\mathcal{V})$ an associative unital graded $\mathbb{F}$-algebra.
\item The symmetric product is well-defined and bilinear.
\item The symmetric product makes $\operatorname{Sym}(\mathcal{V})$ an associative unital graded $\mathbb{F}$-algebra.
\item The exterior product is well-defined and bilinear.
\item The exterior product makes $\bigwedge\mathcal{V}$ an associative unital graded $\mathbb{F}$-algebra.
\end{enumerate}
\begin{proof}
\begin{enumerate}[label=(\roman*)]
\setlength{\itemsep}{0pt}
\item Tensor concatenation is bilinear and associative under the canonical identifications of iterated tensor products.\\
Its unit is $1\in T_0$, and $(\forall\,p,q\in\mathbb{Z}_{\geq0})\,T_pT_q\subseteq T_{p+q}$ by construction.
\item Let $p,q\in\mathbb{Z}_{\geq0}$. Extend $\sigma\in G_p$ to $\widehat{\sigma}\in G_{p+q}$ by fixing the last $q$ indices.\\
Extend $\tau\in G_q$ to $\widehat{\tau}\in G_{p+q}$ by fixing the first $p$ indices.\\
For $\mathbf{x}\in T_p$ and $\mathbf{y}\in T_q$, $(\lambda_{\sigma}\mathbf{x}-\mathbf{x})\otimes\mathbf{y}=\lambda_{\widehat{\sigma}}(\mathbf{x}\otimes\mathbf{y})-\mathbf{x}\otimes\mathbf{y}\in R_{\vee,p+q}$.\\
Similarly, $\mathbf{x}\otimes(\lambda_{\tau}\mathbf{y}-\mathbf{y})=\lambda_{\widehat{\tau}}(\mathbf{x}\otimes\mathbf{y})-\mathbf{x}\otimes\mathbf{y}\in R_{\vee,p+q}$.\\
Hence $R_{\vee,p}\otimes T_q+T_p\otimes R_{\vee,q}\subseteq R_{\vee,p+q}$.\\
The quotient formula is therefore independent of representatives, and bilinearity descends from concatenation.
\item Let $p,q,r\in\mathbb{Z}_{\geq0}$ and $(\mathbf{x},\mathbf{y},\mathbf{z})\in T_p\times T_q\times T_r$. Set $\mathbf{a}\defeq q_{\vee,p}(\mathbf{x})$, $\mathbf{b}\defeq q_{\vee,q}(\mathbf{y})$, and $\mathbf{c}\defeq q_{\vee,r}(\mathbf{z})$.\\
Then $(\mathbf{a}\vee\mathbf{b})\vee\mathbf{c}=q_{\vee,p+q+r}((\mathbf{x}\otimes\mathbf{y})\otimes\mathbf{z})=q_{\vee,p+q+r}(\mathbf{x}\otimes(\mathbf{y}\otimes\mathbf{z}))=\mathbf{a}\vee(\mathbf{b}\vee\mathbf{c})$.\\
The identities $1\vee\mathbf{a}=\mathbf{a}=\mathbf{a}\vee1$ hold by the tensor unit identifications.\\
For $\mathbf{a}\in\operatorname{Sym}^{p}(\mathcal{V})$ and $\mathbf{b}\in\operatorname{Sym}^{q}(\mathcal{V})$, $\mathbf{a}\vee\mathbf{b}\in\operatorname{Sym}^{p+q}(\mathcal{V})$ by definition.
\item Let $p,q\in\mathbb{Z}_{\geq0}$. For $(\mathbf{v},\mathbf{w})\in\mathcal{V}^{\times p}\times\mathcal{V}^{\times q}$, $t_p(\mathbf{v})\otimes t_q(\mathbf{w})\in R_{\wedge,p+q}$ if $\mathbf{v}\in\Delta_p(\mathcal{V})$ or $\mathbf{w}\in\Delta_q(\mathcal{V})$. For $k=0$, $T_k$ is spanned by $t_k(())$.\\
For $k>0$, decomposable tensors span $T_k$ by Theorem \ref{thm:iterated-tensor-products-basic-laws}(i)(b). Hence $R_{\wedge,p}\otimes T_q+T_p\otimes R_{\wedge,q}\subseteq R_{\wedge,p+q}$.\\
The quotient formula is therefore independent of representatives, and bilinearity descends from concatenation.
\item Let $p,q,r\in\mathbb{Z}_{\geq0}$ and $(\mathbf{x},\mathbf{y},\mathbf{z})\in T_p\times T_q\times T_r$ Set $\mathbf{a}\defeq q_{\wedge,p}(\mathbf{x})$, $\mathbf{b}\defeq q_{\wedge,q}(\mathbf{y})$, and $\mathbf{c}\defeq q_{\wedge,r}(\mathbf{z})$.\\
Then $(\mathbf{a}\wedge\mathbf{b})\wedge\mathbf{c}=q_{\wedge,p+q+r}((\mathbf{x}\otimes\mathbf{y})\otimes\mathbf{z})=q_{\wedge,p+q+r}(\mathbf{x}\otimes(\mathbf{y}\otimes\mathbf{z}))=\mathbf{a}\wedge(\mathbf{b}\wedge\mathbf{c})$.\\
The identities $1\wedge\mathbf{a}=\mathbf{a}=\mathbf{a}\wedge1$ hold by the tensor unit identifications. For $\mathbf{a}\in\bigwedge^{p}\mathcal{V}$ and $\mathbf{b}\in\bigwedge^{q}\mathcal{V}$, $\mathbf{a}\wedge\mathbf{b}\in\bigwedge^{p+q}\mathcal{V}$ by definition.
\end{enumerate}
\hlt{Structural Commentary}: Block-relation stability descends concatenation; tensor identities transfer the algebra laws.\\
\hlt{Key Insight}: Adjoining tensor factors preserves both relation subspaces.
\end{proof}
\end{proposition}
